% ============================================================
% Paper 94 (DE): Hartshornes Vermutungen in der Algebraischen Geometrie
% Autor: Michael Fuhrmann
% Projekt: Specialist Mathematics Project
% Build: 169
% Datum: 2026-03-12
% ============================================================
\documentclass[12pt,a4paper]{amsart}

\usepackage{amsmath,amssymb,amsthm,mathtools}
\usepackage{enumitem}
\usepackage{booktabs}
\usepackage{array}
\usepackage[hidelinks]{hyperref}
\usepackage[utf8]{inputenc}
\usepackage[T1]{fontenc}
\usepackage[ngerman]{babel}
\usepackage{geometry}
\geometry{margin=2.5cm}

% ---- Theorem-Umgebungen -------------------------------------
\newtheorem{theorem}{Theorem}[section]
\newtheorem{lemma}[theorem]{Lemma}
\newtheorem{corollary}[theorem]{Korollar}
\newtheorem{proposition}[theorem]{Proposition}
\theoremstyle{definition}
\newtheorem{definition}[theorem]{Definition}
\newtheorem{example}[theorem]{Beispiel}
\theoremstyle{remark}
\newtheorem{remark}[theorem]{Bemerkung}
\newtheorem{conjecture}[theorem]{Vermutung}
% Deutsche Aliasse
\newtheorem{satz}[theorem]{Satz}
\newtheorem{korollar}[theorem]{Korollar}
\newtheorem{vermutung}[theorem]{Vermutung}
\newtheorem{bemerkung}[theorem]{Bemerkung}
\newtheorem{definition_de}[theorem]{Definition}

% ---- Operatoren ---------------------------------------------
\DeclareMathOperator{\codim}{codim}
\DeclareMathOperator{\rank}{rank}
\DeclareMathOperator{\Ext}{Ext}
\DeclareMathOperator{\Hom}{Hom}
\DeclareMathOperator{\Pic}{Pic}
\DeclareMathOperator{\Spec}{Spec}
\newcommand{\PP}{\mathbb{P}}
\newcommand{\ZZ}{\mathbb{Z}}
\newcommand{\CC}{\mathbb{C}}
\newcommand{\OO}{\mathcal{O}}

% ---- Zusammenfassung ----------------------------------------
\begin{abstract}
Wir geben einen Überblick über Hartshornes Vermutungen zu vollständigen
Durchschnitten und Vektorbündeln im projektiven Raum.
Der Quillen--Suslin-Satz (1976), der Serres Vermutung über projektive
Moduln über Polynomringen bewies, wird detailliert dargestellt.
Wir besprechen den Barth--Lefschetz-Satz, das Horrocks--Mumford-Bündel
vom Rang 2 auf $\PP^4$, die Aufspaltungsvermutung für niedrigrangige
Vektorbündel auf $\PP^n$ sowie den
Hartshorne--Lichtenbaum-Verschwindungssatz.
Hartshornes Kodimension-2-Vermutung---dass jede glatte Untervarietät der
Kodimension 2 in $\PP^n$ für $n \ge 7$ ein vollständiger Durchschnitt ist---
bleibt offen.
\end{abstract}

\title{Hartshornes Vermutungen in der Algebraischen Geometrie}

\author{Michael Fuhrmann}
\address{Specialist Mathematics Project, Build 169}
\email{specialist-maths@project.local}
\date{2026-03-12}

\maketitle
\tableofcontents

% =============================================================
\section{Einleitung}
% =============================================================

Der projektive Raum $\PP^n_\CC$ ist die einfachste kompakte algebraische
Varietät, doch das Verständnis, welche Untervarietäten darin existieren können,
ist tiefgründig.  Eine der elegantesten Fragen lautet: Welche glatten
Untervarietäten sind vollständige Durchschnitte?

Eine glatte Untervarietät $Y \subseteq \PP^n$ der Kodimension $c$ ist ein
\emph{vollständiger Durchschnitt}, wenn ihr Ideal $I(Y)$ von genau $c$
homogenen Polynomen erzeugt wird.  In seiner Arbeit von 1974 \cite{Hartshorne1974}
und seinem Buch \cite{Hartshorne1977} formulierte Hartshorne mehrere Vermutungen,
die nahelegen, dass in hochdimensionalen projektiven Räumen glatte
Untervarietäten kleiner Kodimension vollständige Durchschnitte sein müssen.

Gleichzeitig wurde Serres Vermutung---projektive Moduln über Polynomringen
$k[x_1,\ldots,x_n]$ sind frei---unabhängig von Quillen \cite{Quillen1976}
und Suslin \cite{Suslin1976} im Jahr 1976 bewiesen.

% =============================================================
\section{Vollständige Durchschnitte im projektiven Raum}
% =============================================================

\begin{definition_de}[Vollständiger Durchschnitt]
Ein abgeschlossenes Unterschema $Y \subseteq \PP^n_k$ der Kodimension $c$
ist ein \emph{(schematheoretischer) vollständiger Durchschnitt}, wenn die
Idealgarbe $\mathcal{I}_Y$ von genau $c$ homogenen Polynomen erzeugt wird:
$\mathcal{I}_Y = (f_1,\ldots,f_c)$.
\end{definition_de}

\begin{example}
\begin{enumerate}[label=(\roman*)]
  \item Hyperflächen in $\PP^n$ (Kodimension 1) sind immer vollständige
        Durchschnitte.
  \item Die torische kubische Kurve $C = \{[s^3:s^2t:st^2:t^3]\} \subseteq \PP^3$
        vom Grad 3 ist \emph{kein} vollständiger Durchschnitt:
        ihr Ideal benötigt 3 Erzeuger.
  \item Die Veronese-Fläche $V \subseteq \PP^5$ (Kodimension 3) ist kein
        vollständiger Durchschnitt.
\end{enumerate}
\end{example}

% =============================================================
\section{Der Barth--Lefschetz-Satz}
% =============================================================

\begin{satz}[Barth, 1970 \cite{Barth1970}]
Sei $Y \subseteq \PP^n_\CC$ eine glatte Untervarietät der Kodimension $c$.
Dann ist die Einschränkungsabbildung
\[
  H^k(\PP^n, \ZZ) \longrightarrow H^k(Y, \ZZ)
\]
ein Isomorphismus für $k < n - 2c + 1$ und injektiv für $k = n - 2c + 1$.
\end{satz}

\begin{bemerkung}
Für einen vollständigen Durchschnitt der Kodimension $c$ liefert der
Lefschetz-Hyperebenen-Satz Isomorphismen im weiteren Bereich $k \le n - c - 1$.
Barths Satz ist das Analogon ohne die Vollständiger-Durchschnitt-Hypothese,
aber mit einer engeren Schranke.
\end{bemerkung}

\begin{korollar}
Ist $Y \subseteq \PP^n$ glatt mit $\codim Y = c$ und $n \ge 2c$, so gilt
$\pi_1(Y) = 1$ und $\Pic(Y) \cong \ZZ$.
\end{korollar}

% =============================================================
\section{Hartshornes Kodimension-2-Vermutung}
% =============================================================

\begin{vermutung}[Hartshorne, 1974 \cite{Hartshorne1974} -- OFFEN]
\label{verm:hartshorne}
Sei $Y \subseteq \PP^n_\CC$ eine glatte Untervarietät der Kodimension 2.
Falls $n \ge 7$, dann ist $Y$ ein vollständiger Durchschnitt.
\end{vermutung}

\begin{bemerkung}
Die Schranke $n \ge 7$ ist wesentlich:
\begin{itemize}
  \item $n = 3$: Viele nicht-vollständige Durchschnitte (z.\,B.\ torische Kubik).
  \item $n = 4$: Die Horrocks--Mumford-Fläche ist ein Nicht-Durchschnitt.
  \item $n = 5, 6$: Vermutung ebenfalls offen.
  \item $n \ge 7$: Hartshornes Vermutung (offen, keine Gegenbeispiele bekannt).
\end{itemize}
\end{bemerkung}

\begin{satz}[Teilresultat, Hartshorne--Ogus \cite{HartshorneOgus1974}]
Ist $Y \subseteq \PP^n$ glatt der Kodimension 2 und $n \ge 2\dim Y + 1$,
dann zerfällt das Normalenbündel $N_{Y/\PP^n}$ als direkte Summe von
Geradenbündeln.  In diesem Bereich hat $Y$ die kohomologischen Eigenschaften
eines vollständigen Durchschnitts.
\end{satz}

% =============================================================
\section{Die Serre-Korrespondenz und Vektorbündel}
% =============================================================

\begin{satz}[Serre-Korrespondenz]
Sei $Y \subseteq \PP^n$ eine lokal Cohen--Macaulay Untervarietät der Kodimension 2.
Dann gibt es ein Rang-2-Vektorbündel $\mathcal{E}$ auf $\PP^n$ und einen
Schnitt $s \in H^0(\PP^n, \mathcal{E}(t))$ für ein $t$, so dass $Y = Z(s)$
der Nullstellenort von $s$ ist.
\end{satz}

\begin{korollar}
Hartshornes Vermutung ist äquivalent zu: Jedes Rang-2-Bündel auf $\PP^n$
für $n \ge 7$ ist eine direkte Summe von Geradenbündeln $\OO(a) \oplus \OO(b)$.
\end{korollar}

% =============================================================
\section{Die Aufspaltungsvermutung für Vektorbündel}
% =============================================================

\begin{vermutung}[Aufspaltungsvermutung -- OFFEN]
\label{verm:aufspaltung}
Jedes algebraische Vektorbündel vom Rang $r \le \lfloor n/2 \rfloor$
auf $\PP^n_\CC$ zerfällt als direkte Summe von Geradenbündeln.
\end{vermutung}

\begin{satz}[Horrocks-Kriterium \cite{Horrocks1964}]
Ein Vektorbündel $\mathcal{E}$ auf $\PP^n$ zerfällt als direkte Summe von
Geradenbündeln genau dann, wenn $H^i(\PP^n, \mathcal{E}(k)) = 0$ für alle
$k \in \ZZ$ und $1 \le i \le n-1$ gilt.
\end{satz}

\begin{center}
\begin{tabular}{clll}
\toprule
$n$ & Rang $r$ & Status & Quelle \\
\midrule
$\PP^1$ & beliebig & Immer Aufspaltung & Klassisch \\
$\PP^2$ & beliebig & Nicht-Aufspaltung möglich ($r \ge 2$) & \\
$\PP^3$ & $r=2$ & Nicht-Aufspaltung möglich & \\
$\PP^4$ & $r=2$ & Horrocks--Mumford existiert & \cite{HorrocksMumford1973} \\
$\PP^n$, $n \ge 6$ & $r=2$ & Aufspaltungsvermutung (offen) & \cite{Hartshorne1974} \\
\bottomrule
\end{tabular}
\end{center}

% =============================================================
\section{Das Horrocks--Mumford-Bündel}
% =============================================================

\begin{satz}[Horrocks--Mumford, 1973 \cite{HorrocksMumford1973}]
Es gibt ein unzerlegbares Rang-2-Vektorbündel $\mathcal{F}_\mathrm{HM}$
auf $\PP^4_\CC$ (das \emph{Horrocks--Mumford-Bündel}) mit Chern-Klassen
$c_1 = 5$, $c_2 = 10$.  Die Nullstellenorte seiner Schnitte sind abelsche
Flächen vom Grad 10 in $\PP^4$.
\end{satz}

\begin{bemerkung}
Das Horrocks--Mumford-Bündel ist das einzige bekannte unzerlegbare
Rang-2-Bündel auf $\PP^n$ für $n \ge 4$.  Seine Existenz zeigt, dass die
Aufspaltungsvermutung für $n = 4$ falsch ist.
Die Frage für $n \ge 5$ ist offen.
\end{bemerkung}

Die Konstruktion von $\mathcal{F}_\mathrm{HM}$ nutzt die Wirkung der
Heisenberg-Gruppe $H_5$ auf $\PP^4$ und die Theorie der Theta-Funktionen.

% =============================================================
\section{Serres Vermutung: Der Quillen--Suslin-Satz}
% =============================================================

\begin{satz}[Quillen--Suslin, 1976 \cite{Quillen1976,Suslin1976}]
\label{satz:quillen-suslin}
Sei $k$ ein Körper (oder ein Hauptidealbereich).  Jeder endlich erzeugte
projektive Modul über dem Polynomring $k[x_1,\ldots,x_n]$ ist frei.
\end{satz}

\begin{bemerkung}
Serre formulierte dies 1955 als Vermutung \cite{Serre1955}.
Geometrisch bedeutet es: Jedes algebraische Vektorbündel auf dem affinen
Raum $\mathbb{A}^n_k$ ist trivial.
\end{bemerkung}

\begin{proof}[Beweisidee (Quillens Ansatz)]
\textbf{Schritt 1.}  Nach einem Satz von Horrocks (1964) ist ein projektiver
Modul $P$ über $k[x_1,\ldots,x_n]$ genau dann frei, wenn
$P \otimes_{k[x_n]} k[x_n]_{(f)}$ für ein einziges monisches Polynom
$f \in k[x_n]$ frei ist.

\textbf{Schritt 2.}  Quillens \emph{Patch-Lemma}: Ist $M$ ein endlich
präsentierter $A[t]$-Modul, der für alle maximalen Ideale
$\mathfrak{m} \subset A$ von $A_\mathfrak{m}$ ausgedehnt ist, dann ist $M$
von $A$ ausgedehnt.

\textbf{Schritt 3.}  Induktion über $n$: Für $n=1$ sind projektive Moduln
über einem Hauptidealbereich frei.  Das Patch-Lemma reduziert auf den
$n{-}1$-Fall. $\square$
\end{proof}

% =============================================================
\section{Der Hartshorne--Lichtenbaum-Verschwindungssatz}
% =============================================================

\begin{satz}[Hartshorne--Lichtenbaum-Verschwindungssatz \cite{Hartshorne1968}]
Sei $X$ eine glatte projektive Varietät der Dimension $n$ über einem Körper $k$,
und $Y \subseteq X$ ein abgeschlossenes Unterschema.  Dann gilt
\[
  H^n_Y(X, \OO_X) = 0,
\]
sofern $Y$ keine irreduzible Komponente von $X$ enthält und
$\codim_X(Y) \ge 1$.  Allgemeiner gilt für ein amples Geradenbündel $\mathcal{L}$:
\[
  H^i(X, \mathcal{L}^{-1}) = 0 \quad \text{für } i < n.
\]
\end{satz}

\begin{bemerkung}
Dies verallgemeinert den Kodaira-Verschwindungssatz und ist ein Grundpfeiler
der modernen algebraischen Geometrie.
\end{bemerkung}

% =============================================================
\section{Algebraische K-Theorie und Vektorbündel}
% =============================================================

Serres Vermutung und die Aufspaltungsvermutung fügen sich in den Rahmen der
algebraischen K-Theorie ein.

\begin{definition_de}[Grothendieck-Gruppe]
Für ein Schema $X$ ist die \emph{Grothendieck-Gruppe} $K_0(X)$ die freie
abelsche Gruppe auf Isomorphieklassen kohärenter Garben auf $X$, modulo der
Relation $[\mathcal{F}] = [\mathcal{F}'] + [\mathcal{F}'']$ für jede kurze
exakte Sequenz $0 \to \mathcal{F}' \to \mathcal{F} \to \mathcal{F}'' \to 0$.
\end{definition_de}

\begin{satz}[Quillen, 1973 \cite{Quillen1973}]
$K_i(\PP^n_k) \cong K_i(k)^{n+1}$ für alle $i \ge 0$
(projektive Bündelformel in höherer K-Theorie).
\end{satz}

Für den affinen Raum: $K_i(\mathbb{A}^n_k) \cong K_i(k)$ für alle $i$
(Homotopie-Invarianz), was die K-theoretische Form des Quillen--Suslin-Satzes ist.

% =============================================================
\section{Aktueller Stand und verwandte Resultate}
% =============================================================

\begin{satz}[Faltings, 1981 \cite{Faltings1981}]
Ist $Y \subseteq \PP^n$ ein glatter vollständiger Durchschnitt der
Kodimension $c$ und $n > 2c$, dann ist $Y$ ein schematheoretischer
vollständiger Durchschnitt.
\end{satz}

\begin{satz}[Evans--Griffith, 1981 \cite{EvansGriffith1985}]
Sei $\mathcal{E}$ ein Vektorbündel vom Rang $r$ auf $\PP^n_k$.
Wenn $\mathcal{E}$ auf jedem $(r-1)$-dimensionalen projektiven Unterraum
zerfällt, dann zerfällt $\mathcal{E}$.
\end{satz}

% =============================================================
\section{Offene Probleme}
% =============================================================

\begin{vermutung}[Hartshornes Kodimension-2-Vermutung, = Vermutung~\ref{verm:hartshorne} -- OFFEN]
Jede glatte Untervarietät der Kodimension 2 in $\PP^n_\CC$ für $n \ge 7$
ist ein vollständiger Durchschnitt.
\end{vermutung}

\begin{vermutung}[Aufspaltungsvermutung für Rang-2-Bündel, = Vermutung~\ref{verm:aufspaltung} -- OFFEN]
Jedes algebraische Rang-2-Vektorbündel auf $\PP^n_\CC$ für $n \ge 6$
zerfällt als $\OO(a) \oplus \OO(b)$.
\end{vermutung}

\begin{vermutung}[Mengentheoretische vollständige Durchschnitte -- teilweise offen]
Jede glatte Untervarietät $Y \subseteq \PP^n$ ist ein \emph{mengentheoretischer}
vollständiger Durchschnitt (d.\,h.\ als Menge durch $\codim Y$ Gleichungen
definiert).  Dies ist in Charakteristik $p > 0$ bekannt (Kempf), aber in
Charakteristik 0 für Kurven in $\PP^3$ offen.
\end{vermutung}

\begin{vermutung}[Hartshornes allgemeine Vermutung \cite{Hartshorne1974}]
Ist $Y$ eine glatte Untervarietät von $\PP^n$ mit $\dim Y > \tfrac{2n}{3}$,
dann ist $Y$ ein vollständiger Durchschnitt.
\end{vermutung}

% =============================================================
\section{Fazit}
% =============================================================

Hartshornes Vermutungen haben eine Generation von Forschung an der
Schnittstelle von algebraischer Geometrie, kommutativer Algebra und Topologie
angetrieben.  Der Quillen--Suslin-Satz hat den affinen Fall vollständig
gelöst; der projektive Fall durch die Aufspaltungsvermutung und Hartshornes
Kodimension-2-Vermutung bleibt weitgehend offen.

Das Horrocks--Mumford-Bündel demonstriert die Subtilität des projektiven
Falls: Unzerlegbare Rang-2-Bündel existieren auf $\PP^4$, aber keines ist
auf $\PP^n$ für $n \ge 5$ bekannt.  Ob dies eine echte geometrische Schranke
oder eine Lücke in unseren Methoden widerspiegelt, ist unbekannt.

\begin{thebibliography}{99}

\bibitem{Hartshorne1974}
R.~Hartshorne,
\textit{Varieties of small codimension in projective space},
Bull. Amer. Math. Soc. \textbf{80} (1974), 1017--1032.

\bibitem{Hartshorne1977}
R.~Hartshorne,
\textit{Algebraic Geometry},
Graduate Texts in Mathematics \textbf{52},
Springer, New York, 1977.

\bibitem{Quillen1976}
D.~Quillen,
\textit{Projective modules over polynomial rings},
Invent. Math. \textbf{36} (1976), 167--171.

\bibitem{Suslin1976}
A.~A. Suslin,
\textit{Projektive Moduln über einem Polynomring},
Dokl. Akad. Nauk SSSR \textbf{226} (1976), 549--551.

\bibitem{Barth1970}
W.~Barth,
\textit{Transplanting cohomology classes in complex-projective space},
Amer. J. Math. \textbf{92} (1970), 951--967.

\bibitem{HorrocksMumford1973}
G.~Horrocks und D.~Mumford,
\textit{A rank 2 vector bundle on $\mathbf{P}^4$ with 15{,}000 symmetries},
Topology \textbf{12} (1973), 63--81.

\bibitem{Horrocks1964}
G.~Horrocks,
\textit{Vector bundles on the punctured spectrum of a local ring},
Proc. London Math. Soc. (3) \textbf{14} (1964), 689--713.

\bibitem{Serre1955}
J.-P. Serre,
\textit{Faisceaux algébriques cohérents},
Ann. of Math. (2) \textbf{61} (1955), 197--278.

\bibitem{Hartshorne1968}
R.~Hartshorne,
\textit{Cohomological dimension of algebraic varieties},
Ann. of Math. (2) \textbf{88} (1968), 403--450.

\bibitem{HartshorneOgus1974}
R.~Hartshorne und A.~Ogus,
\textit{On the factoriality of local rings of small embedding codimension},
Comm. Algebra \textbf{1} (1974), 415--437.

\bibitem{Quillen1973}
D.~Quillen,
\textit{Higher algebraic K-theory I},
in: \textit{Algebraic K-theory I}, Lecture Notes in Math. \textbf{341},
Springer, 1973, 85--147.

\bibitem{Faltings1981}
G.~Faltings,
\textit{Ein Kriterium für vollständige Durchschnitte},
Invent. Math. \textbf{62} (1981), 393--401.

\bibitem{EvansGriffith1985}
E.~G. Evans und P.~Griffith,
\textit{The syzygy problem},
Ann. of Math. (2) \textbf{114} (1981), 323--333.

\bibitem{Zak1993}
F.~L. Zak,
\textit{Tangents and Secants of Algebraic Varieties},
Translations of Math. Monographs \textbf{127},
American Mathematical Society, 1993.

\end{thebibliography}

\end{document}
